\chapter{Toric CY3 and Critical CoHAs}
\label{ch:toric-coha}

This chapter develops the second main class of examples: toric Calabi--Yau threefolds, whose critical cohomological Hall algebras provide the positive half of affine super Yangians.

\section{Toric CY3 threefolds and their quivers}
\label{sec:toric-cy3-quivers}

A toric CY3 $X_\Sigma$ is determined by a fan $\Sigma$ in $\mathbb{Z}^3$ satisfying the CY condition (all generators coplanar).  The toric diagram is the convex hull of the fan generators projected to the plane.

\begin{example}[$\mathbb{C}^3$]
\label{ex:c3}
The simplest toric CY3.  Toric diagram: a single triangle.  Quiver: the Jordan quiver (one vertex, one loop).
\end{example}

\begin{example}[Resolved conifold]
\label{ex:conifold}
Toric diagram: two triangles sharing an edge.  Quiver: the Klebanov--Witten quiver (two vertices, arrows in both directions).
\end{example}

\section{The critical cohomological Hall algebra}
\label{sec:critical-coha}

\begin{definition}[Critical CoHA]
\label{def:critical-coha}
For a quiver with potential $(Q, W)$ coming from a toric CY3 $X$, the \emph{critical CoHA} is the $\mathbb{Z}$-graded associative algebra
\[
  \mathcal{H}(Q, W) = \bigoplus_{\mathbf{d} \in \mathbb{Z}_{\geq 0}^{Q_0}} H^{\mathrm{BM}}_*(\mathrm{Crit}(W_\mathbf{d}), \phi_{W_\mathbf{d}})
\]
where $\phi_{W_\mathbf{d}}$ denotes the vanishing cycle sheaf and $W_\mathbf{d}$ is the potential restricted to the representation space of dimension vector $\mathbf{d}$.  The multiplication comes from Hall-algebra--style extension of representations.
\end{definition}

\section{$\mathbb{C}^3$ and the affine Yangian $Y(\widehat{\mathfrak{gl}}_1)$}
\label{sec:c3-yangian}

\begin{theorem}[Schiffmann--Vasserot]
\label{thm:sv-c3}
The critical CoHA of $(\mathbb{C}^3, 0)$ (Jordan quiver, cubic potential $W = \mathrm{tr}(XYZ - XZY)$) is isomorphic to the positive half of the affine Yangian of $\mathfrak{gl}_1$:
\[
  \mathcal{H}(\mathbb{C}^3) \simeq Y^+(\widehat{\mathfrak{gl}}_1).
\]
\end{theorem}

The full affine Yangian $Y(\widehat{\mathfrak{gl}}_1)$ is isomorphic to the $\mathcal{W}_{1+\infty}$ algebra at the self-dual level --- already a key object in the standard landscape of Volume~I, where MC4 is solved by weight stabilization.  The Yangian $R$-matrix is the DK-0 shadow.

\section{General toric CY3 and affine super Yangians}
\label{sec:general-toric-yangian}

\begin{theorem}[Rapcak--Soibelman--Yang--Zhao]
\label{thm:rsyz}
For a toric CY3 $X$ without compact $4$-cycles, the critical CoHA $\mathcal{H}(Q_X, W_X)$ is isomorphic to the positive half of the affine super Yangian $Y(\widehat{\mathfrak{g}}_{Q_X})$ associated to the toric quiver.
\end{theorem}

The toric diagram determines the quiver, hence the super Lie algebra $\mathfrak{g}_{Q_X}$, hence the affine super Yangian.  The toric fan IS the root datum of the quantum vertex chiral group $G(X)$.

\section{The CoHA as $E_1$-sector}
\label{sec:coha-e1}

The critical CoHA is an associative algebra --- an $E_1$-algebra.  In the language of this work:
\begin{itemize}
  \item The CoHA = the positive half of $G(X)$ = the $E_1$-chiral sector (the ordered part).
  \item The full quantum vertex chiral group $G(X)$ is $E_2$ (braided).
  \item The braiding (the passage from $E_1$ to $E_2$) is the quantum group $R$-matrix of the affine super Yangian.
  \item The representation category $\Rep^{E_2}(G(X))$ is braided monoidal, with braiding from the Yangian $R$-matrix.
\end{itemize}

This is the ``tree-level'' CY3 combinatorics: the CoHA captures the genus-0 curve counting (DT invariants as intersection numbers on Hilbert schemes), and the E_1 structure encodes the associative (ordered) OPE.

\section{Root datum from toric geometry}
\label{sec:toric-root-datum}

The generalized root datum $\mathcal{R}(X_\Sigma)$ of a toric CY3:
\begin{enumerate}[label=(\roman*)]
  \item \textbf{Lattice}: $\Lambda(X) = K_0(\mathrm{Coh}(X))$ with the Euler form.
  \item \textbf{Real roots}: vertices of the toric diagram (compact divisors).
  \item \textbf{Imaginary roots}: dimension vectors of quiver representations with nonzero DT invariant.
  \item \textbf{Root multiplicities}: $\mathrm{mult}(\mathbf{d}) = \mathrm{DT}_\mathbf{d}(X)$ (Donaldson--Thomas invariants).
  \item \textbf{Weyl group}: symmetries of the toric diagram.
\end{enumerate}
